\chapter{Investigando a Mente}

\lettrine{A}{raiz do sofrimento} é o que chamamos \emph{avijjā} -- o não saber, ou a ignorância de
como realmente as coisas são. Esta ignorância básica é aquela de não compreender
a nossa verdadeira natureza. Sofremos por causa de maneiras de ver e opiniões, e
por causa de hábitos e condições que não compreendemos. Vivemos as nossas vidas
num estado de ignorância, não compreendendo a forma como as coisas são.

Se nos ouvirmos muito, por vezes podemos ouvir declarações tais como: ``Eu devia
fazer isto, eu não devia fazer isto, eu devia ser assim, eu não devia ser
assim'', ou que o mundo seria diferente do que é; os nossos pais deveriam ser
desta ou daquela maneira e não deviam ser como são. Então, temos este tempo
verbal em particular a ressoar na nossa cabeça porque temos uma ideia do que
deveria ou não deveria ser. Na meditação, escutem aquela opinião dentro de vós
do que deveria e não deveria ser -- simplesmente escutem-na.

A nossa tendência é tentar tornarmo-nos algo; então estabelecemos um objectivo,
criamos um ideal do que gostaríamos de nos tornar. Talvez pensemos que a
sociedade deve ser diferente do que é. As pessoas deviam ser amáveis, generosas,
altruístas, compreensivas e amorosas; devia haver fraternidade e as pessoas não
deviam ser egoístas. O governo devia ter líderes sábios, o mundo devia estar em
paz e assim por diante. Mas o mundo está como está neste momento e as coisas são
como são. Quando não compreendemos isto, então debatemo-nos. Assim, ouçam
interiormente o choro constante. ``Eu sou assim, eu não sou assim'', e penetrem
neste ``eu sou, eu não sou'' com consciência.

Temos tendência para reagir e assumir por garantido que todo o ``eu sou'' e
``não sou'' é a verdade. Criamo-nos como uma personalidade e apegamo-nos às
nossas memórias. Lembramo-nos das coisas que aprendemos e do que fizemos --
geralmente as coisas mais extremas; temos tendência para esquecer as coisas mais
comuns. % REVISÃO: "Pessoas que viveram vidas muito egoístas, têm de beber" — a vírgula separa indevidamente o sujeito do verbo; deve ser removida.
Portanto, se fizermos coisas indelicadas, cruéis e tolas, então temos
memórias desagradáveis nas nossas vidas; sentimo-nos envergonhados ou culpados.
Se fizermos coisas boas ou caridosas, então temos boas memórias nas nossas
vidas. Quando começarmos a reflectir sobre esta questão, então teremos mais
cuidado com o que fazemos e dizemos; se tivermos vivido a vida
inconscientemente, agindo impulsivamente por desejo de gratificação imediata ou
por uma intenção de magoar, causar desarmonia ou explorar os outros, as nossas
mentes estarão cheias de memórias muito desagradáveis. Pessoas que viveram vidas
muito egoístas, têm de beber muito ou tomar drogas para se manterem
constantemente ocupadas, de forma a não terem que olhar para as memórias que
surgem na mente.

No processo de despertar na meditação, estamos a trazer consciência para as
condições da mente aqui e agora só por estarmos conscientes desta noção de ``eu
sou, não sou''. Contemplemos os sentimentos de dor ou prazer -- e quaisquer
memórias, pensamentos e opiniões -- como impermanentes, \emph{anicca}. A
característica do transitório é comum a todas as condições. Quantos de vós
passaram o dia a investigar isto de todas as formas possíveis enquanto sentados,
de pé ou deitados? Investiguem o que vêem através dos olhos, o que ouvem através
dos ouvidos, o que saboreiam com a língua, cheiram com o nariz, sentem e
experimentam com o corpo e pensam com a mente.

O pensamento ``eu sou'' é uma condição impermanente. O pensamento ``não sou'' é
uma condição impermanente. Pensamentos, memórias, consciência do pensamento, o
próprio corpo, as nossas emoções -- todas as condições mudam. Na prática da
meditação, temos de assumir bastante seriedade, bravura e coragem para realmente
investigarmos, ousar olharmos até para as condições mais desagradáveis da vida,
em vez de procurarmos uma fuga na tranquilidade ou de esquecermos tudo. Em
\emph{Vipassana}, a prática é a de olhar para dentro do sofrimento; é um
confronto connosco mesmos, com o que pensamos de nós mesmos, com as nossas
memórias e emoções, agradáveis, desagradáveis ou indiferentes. Por outras
palavras, quando estas coisas surgem e estamos conscientes do sofrimento, em vez
de o rejeitarmos, reprimirmos ou ignorá-lo, aproveitamos a oportunidade para o
examinar.

Então o sofrimento é o nosso professor. Ensina-nos e por isso temos de aprender
a lição através do estudo do próprio sofrimento. Surpreende-me sempre como
algumas pessoas pensam que nunca sofrem. Pensam: ``Não sofro. Não sei porque é
que os budistas estão sempre a falar de sofrimento. Sinto-me maravilhoso, cheio
de beleza e alegria. Estou sempre tão feliz. Acho a vida uma experiência
fantástica, interessante, fascinante e um prazer interminável''. Estas pessoas
têm tendência para aceitar esse lado da vida e rejeitar o outro, porque
inevitavelmente o que nos encanta desaparece e depois lamentamos.

% REVISÃO: "O nosso desejo ... de deleite, leva-nos" — a vírgula separa o sujeito do verbo; deve ser removida.
O nosso desejo de estar num constante estado de deleite, leva-nos a todo o tipo
de problemas, dificuldades e situações. O sofrimento não existe só por causa de
coisas gravíssimas como ter cancro terminal, ou perder alguém que se ama; o
sofrimento pode ocorrer em torno do que é muito comum, como as quatro posturas,
sentar, ficar de pé, andar, deitar-se. Não há nada de extremo nisso.

Contemplamos a respiração normal e a consciência corrente. De forma a
compreender a existência, contemplamos sentimentos, memórias e pensamentos
comuns em vez de procurarmos agarrar ideias e pensamentos fantásticos para
compreendermos os extremos da existência. Assim não nos estamos a envolver com
especulações sobre o propósito final da vida, Deus, o diabo, o céu, o inferno e
o que acontece quando morremos ou reencarnamos. Na meditação budista, só
observamos o aqui e agora. O nascimento e morte que está a acontecer aqui e
agora é o começo e o fim das coisas mais comuns.

Contemple-se o início. Quando se pensa em nascimento, pensa-se em ``Eu nasci'',
mas esse é o grande nascimento do corpo, que não nos conseguimos lembrar. O
nascimento comum de ``eu'' que vivenciamos na vida diária é ``eu quero, eu não
quero, eu gosto, eu não gosto''. Isto é um nascimento, ou busca de ser feliz.
Contemplamos o comum inferno da nossa própria raiva, a raiva que surge, o calor
do corpo, a aversão, o ódio que sentimos na mente. Contemplamos o paraíso comum
que experimentamos, os estados felizes, a bem-aventurança, a leveza, a beleza no
aqui e agora. Ou apenas o estado mental embotado, o tipo de limbo, nem feliz nem
infeliz, mas embotado, entediado e indiferente. Na meditação budista, observamos
isto dentro de nós mesmos.

Contemplamos o nosso próprio desejo de poder e controlo, de controlar outro
alguém, de nos tornarmos famosos ou de nos tornarmos alguém que está no topo.
Quantos de vós, quando descobrem que alguém é mais talentoso do que vocês,
querem colocá-lo para baixo? Isso é ciúme. O que temos que fazer na nossa
prática de meditação é ver os ciúmes típicos, ou o ódio que podemos sentir por
alguém que pode aproveitar-se de nós ou incomodar-nos; a ganância ou luxúria que
podemos sentir por alguém que nos atrai. A nossa própria mente é como um espelho
que reflecte o universo e nós observamos o reflexo. Antes, assumiríamos esses
reflexos como realidade de modo a ficarmos fascinados, enojados ou indiferentes
a eles. Mas em \emph{Vipassana}, apenas observamos que todas essas reflexões são
condições mutáveis. % REVISÃO: "onde quando como ignorantes" está truncado/agramatical. Reformular, p. ex.: "ao passo que, como ignorantes, temos a tendência para procurar identidade nelas."
Começamos a vê-las como objectos em vez de um eu, onde
quando como ignorantes, temos a tendência para procurar identidade nelas.

% REVISÃO: concordância de tempos verbais — "sentar-se-á ... e experimentará ... até mesmo cheirar odores ... e ouvir sons divinos e pense:" mistura futuro e infinitivo. Uniformizar no futuro: "cheirará ... ouvirá ... e pensará:".
% REVISÃO: "Por que algo assim nunca acontece comigo?" — construção do português do Brasil; em PT-PT preferir "Porque é que algo assim nunca me acontece?".
Portanto, na prática, estamos a olhar para o universo conforme ele está a ser
reflectido nas nossas mentes. Não importa o que os outros experimentem; um
meditador sentar-se-á aqui e experimentará todos os tipos de luzes brilhantes,
cores, imagens fascinantes, Buddhas, seres celestiais, até mesmo cheirar odores
maravilhosos e ouvir sons divinos e pense: ``Que meditação maravilhosa, veio
tanto brilho, `o esplendor' -- um ser divino veio como um anjo esplendoroso,
tocou-me e senti este êxtase. A experiência de êxtase mais maravilhosa de toda a
minha vida\ldots{} Esperei toda a minha vida por esta experiência.'' Enquanto
isso, o próximo está a pensar: ``Por que algo assim nunca acontece comigo?
Sentei-me durante uma hora inteira com dores nas costas, deprimido, com vontade
de fugir, perguntando-me por que raio vim eu a este retiro''. Outra pessoa
poderá dizer: ``Eu não aguento todas aquelas pessoas que têm aquelas ideias
parvas e fantasias. Elas metem-me nojo; provocam-me este terrível ódio e aversão
em mim. Eu odeio a imagem do Buddha sentado na janela. Quero esmagá-la. Eu odeio
o Budismo e a meditação!''

% REVISÃO: "alguém que vê \emph{devas} e experiência luzes brilhantes" e "do êxtase ou experiência entorpecimento" — aqui "experiência" é forma verbal e deve escrever-se "experiencia" (sem acento, do verbo experienciar).
Agora, qual dessas três pessoas é o bom meditador? Compare-se com aquele que vê
\emph{devas} a dançar no céu, aquele que está aborrecido, indiferente e
embotado, ou aquele cheio de ódio e aversão? \emph{Devas} e anjos a dançar nos
reinos celestiais são \emph{anicca}, são impermanentes. O aborrecimento é
\emph{anicca}, impermanente. O ódio e a aversão são \emph{anicca},
impermanentes. Portanto, o bom meditador, aquele que pratica da maneira certa vê
a natureza impermanente destas condições. Quando fala com alguém que vê
\emph{devas} e experiência luzes brilhantes, começa a duvidar da sua própria
prática e pensa: ``Talvez eu não seja capaz de iluminação. Talvez eu não esteja
a meditar como deve ser''. A dúvida é em si impermanente. Tudo o que surge
passa. Portanto, o bom meditador é aquele que vê a natureza impermanente da
bem-aventurança e do êxtase ou experiência entorpecimento, raiva, ódio e aversão
e reflecte sobre a natureza impermanente dessas qualidades quando medita
sentado, a caminhar ou deitado.

Qual é a sua tendência? É muito optimista sobre tudo? ``Eu gosto de todos aqui.
Eu acredito nos ensinamentos do Buddha, eu acredito no Dhamma'' -- esse tipo de
mente é de fé, acredita. E este tipo de mente pode criar e experienciar muito
rapidamente coisas arrebatadoras. Alguns dos agricultores na Tailândia, pessoas
quase sem conhecimento algum sobre as coisas do mundo, que mal sabem ler e
escrever, às vezes podem experienciar estados de felicidade, luzes, \emph{devas}
e tudo mais, e acreditar neles. Quando acreditamos em \emph{devas}, vemo-los.
Quando acreditamos em luzes e reinos celestiais, iremos vê-los. Acreditamos que
Buddha virá salvar-nos, Buddha virá e salvar-nos-á. Aquilo em que acreditamos,
acontece connosco. Se acreditamos em fantasmas, fadas e elfos, veremos que estas
coisas se manifestam para nós. Mas ainda são \emph{anicca}, impermanentes,
transitórios e não-eu.

A maioria das pessoas não acredita em fadas e \emph{devas} e pensam que tais
coisas são uma idiotice. Esse é o tipo de mente negativa, que é desconfiada e
duvidosa e que não acredita em nada. ``Eu não acredito em fadas e \emph{devas}.
Eu não acredito em nada desse género de coisas. Ridículo? Mostre-me uma fada.''
Portanto, a mente muito desconfiada e céptica nunca vê tais coisas.

A fé existe; a dúvida existe. Na prática budista, examinamos a crença e a dúvida
que experienciamos na nossa mente e vemos que estas duas condições estão a
mudar.

Já contemplei a própria dúvida como um sinal. Eu fazia a mim próprio uma
pergunta como, ``Quem sou eu?'' E então ouvia a resposta -- algo como, ``Sumedho
Bhikkhu''. Então pensava: ``Esta não é a resposta; quem és tu realmente?'' Eu
veria o conflito, a habitual reacção para encontrar uma resposta para a
pergunta. Mas eu não aceitaria qualquer resposta conceitual. ``Quem está aqui
sentado? O que é isto? O que é isto aqui? Quem está pensando, afinal? O que é
que pensa?'' Quando um estado de incerteza ou dúvida surgisse, eu apenas olharia
para essa incerteza ou dúvida como um sinal, porque a mente pára aí e fica em
branco, e então surge o vazio.

Achei útil esvaziar a mente perguntando a mim mesmo perguntas irrespondíveis que
fariam com que surgissem dúvidas. A dúvida é uma condição impermanente. A forma,
o conhecido, é impermanente; não saber é impermanente. Alguns dias eu saía e
olhava a natureza e observava-me apenas parado aqui, olhando para o chão. Eu
perguntar-me-ia: ``O solo está separado de mim? Quem é este que vê o chão? \ldots{}
Estas folhas e o solo estão na minha mente ou fora da minha mente? \ldots{} O que é
isso que vê? É o globo ocular? \ldots{} Se eu arrancasse o meu globo ocular, ele
seria separado de mim? \ldots{} Eu ainda veria essas folhas? \ldots{} Eles ainda ali estão
quando não estou a olhar para eles? \ldots{} Quem é que está ciente disso, afinal?''

E o som. Também fiz algumas experiências com som porque os objectos da visão têm
uma certa solidez como esta sala -- parece bastante permanente, sabe, pelo menos
por hoje. Mas o som é verdadeiramente \emph{anicca} -- tente controlar o som.

Investiguei o som perguntando: ``Os meus olhos podem ouvi-lo? Se eu cortar os
meus ouvidos e tímpanos, haverá algum som? Posso ver e ouvir exactamente no
mesmo momento?'' Todos os órgãos dos sentidos e os seus objectos são
impermanentes, mudando as condições. Pense agora: ``Onde está a sua mãe? Onde
está a minha mãe agora?'' Se eu pensar nela no seu apartamento na Califórnia, é
um conceito na mente. Mesmo se eu pensar, ``A Califórnia é ali'', ainda é a
mente a pensar ``ali''. ``Mãe'' é um conceito, não é? Então, onde está a mãe
agora? Ela está na mente: quando a palavra ``mãe'' surge, você ouve a palavra
como um som e ela traz uma imagem mental ou uma memória ou um sentimento de
gosto, antipatia ou indiferença.

% REVISÃO: "Você tem preconceitos e preconceitos sobre raça ou nacionalidade" — palavra repetida; o original distingue dois termos (p. ex. "ideias preconcebidas e preconceitos").
% REVISÃO: "isto é só inerente nessas identidades" — regência incorrecta; "inerente" rege a preposição "a": "inerente a essas identidades".
Todos os conceitos na mente que consideramos realidade devem ser investigados:
saiba o que os conceitos fazem à mente. Observe o prazer que obtém ao pensar em
certos conceitos e o desprazer que outros trazem. Você tem preconceitos e
preconceitos sobre raça ou nacionalidade -- todos estes são conceitos ou
proliferações conceituais. Os homens têm certas atitudes e preconceitos em
relação aos homens: isto é só inerente nessas identidades.

Mas na meditação, ``feminino'' é um conceito e ``masculino'' é um conceito, um
sentimento, uma percepção na mente. Portanto, nesta prática de \emph{Vipassana},
penetramos com discernimento na natureza de todas as condições, grosseiras ou
subtis. A compreensão interior destrói as ilusões que estes conceitos nos dão,
de serem reais. Podem surgir condições; não podemos parar as coisas que nos
afectam na vida -- como o clima, a economia, problemas familiares, o nosso
passado, a nossa oportunidade ou falta de oportunidade. Mas podemos penetrar em
todas estas condições -- que são impermanentes e não-eu. Este é o caminho da
transcendência; transcendendo a condição mortal através da consciência da
condição mortal.

O Buddha é o professor, aquilo que dentro de nós nos lembra de observar a
natureza impermanente de todas as condições e a não assumir nenhuma delas como
realidade. Quando o fazemos, o que acontece? Temos guerras, greves, batalhas e
problemas sem fim que existem no mundo porque os seres ignorantes assumem essas
condições como realidade. Ligam-se ao corpo mortal como uma identidade. Somos
absorvidos nestes vários símbolos e conceitos e, nessa absorção, temos que
nascer e morrer nessas condições. É como apegarmo-nos a algo que está a
mover-se, como a ganância e ser puxado por esse movimento. Então, nascemos e
morremos nesse momento. Mas quando não nos apegamos mais, então estamos a evitar
sofrer do movimento e das limitações da mudança de condições.

Ora, falando assim, as pessoas podem questionar: ``Mas como viver então nesta
sociedade, se tudo é irreal?'' O Buddha fez uma distinção muito clara entre a
realidade convencional e a realidade última. No nível convencional de
existência, usamos convenções que trazem harmonia para nós mesmos e para a
sociedade em que vivemos. Que tipo de convenções trazem harmonia? Bem, coisas
como ser bom, estar atento, não fazer coisas que causam desarmonia, como roubar,
enganar e explorar os outros. Ter respeito e compaixão pelos outros seres, ser
observador, tentar ajudar: todas estas convenções trazem harmonia.

Portanto, no ensino Budista ao nível convencional, vivemos de uma maneira que
apoia a prática do bem e evitamos fazer o mal com o corpo e a fala. Não é como
se estivéssemos a rejeitar o mundo convencional ``Não quero ter nada a ver com
isso porque é uma ilusão'' -- isso é \emph{outra} ilusão. Pensar que o mundo
convencional é uma ilusão é só mais um pensamento.

Na nossa prática, nós vemos que pensamento é pensamento, ``o mundo é uma
ilusão'' é um pensamento. Mas aqui e agora, devemos estar cientes de que tudo o
que temos consciência está a mudar. Vivam conscientemente, ponham esforço e
concentração no que fazem, estejam sentados, a andar, deitados ou a trabalhar.
Independentemente de se ser homem ou mulher, um secretário, secretária, dona de
casa ou trabalhador, executivo, ou seja, o que for, requer esforço e
concentração. Fazer o bem e evitar fazer o mal. Isto é como um budista vive
dentro das formas convencionais da sociedade. Mas não mais iludidos pelo corpo
ou pela sociedade, ou pelas coisas que se passam na sociedade, porque uma pessoa
budista é aquela que examina e investiga o universo ao investigar o seu próprio
corpo e mente.
